% ---------------------------------------------------------------------
\chapter{SW01 – Asymptotische Analyse, O-Notation, Rekursion}
\label{kap:sw01}

\section*{Wochenziele}
\begin{wochenziele}
  \ziel{Arithmetische Folge erkennen}{Eine konstante Differenz zwischen aufeinanderfolgenden Gliedern erkennen.}
  \ziel{Rekursive Form}{\(a_n\) mithilfe des vorherigen Glieds und eines Startwerts definieren.}
  \ziel{Iterative Form}{Das n-te Glied als Summe vorheriger Schritte darstellen.}
  \ziel{Explizite Form}{Eine Formel für \(a_n\) direkt in Abhängigkeit von \(n\) herleiten.}
  \ziel{Startwert bestimmen}{Erkennen, ob die Folge bei \(a_0\) oder \(a_1\) beginnt.}
  \ziel{Differenz \(d\) bestimmen}{Die konstante Differenz zwischen zwei aufeinanderfolgenden Gliedern berechnen.}
  \ziel{Summenzeichen verstehen}{\(\sum\) lesen und die Grenzen der Summe richtig anwenden.}
  \ziel{Arithmetische Reihe berechnen}{\(1+2+\cdots+(n-1)=\frac{n(n-1)}{2}\) anwenden.}
  \ziel{Formel vereinfachen}{Klammern auflösen, Terme zusammenfassen und in Polynom-Normalform bringen.}
  \ziel{Polynome ordnen}{Terme nach absteigenden Potenzen von \(n\) ordnen.}
  \ziel{Zusammenhang zur Laufzeit}{Erkennen, wie Summen aus Schleifen zu \(\Oh{n^2}\) oder \(\Oh{n^3}\) führen können.}
\end{wochenziele}

\section{Kernideen}
% Worum geht es? In eigenen Worten, 3-5 Sätze.

\section{Definitionen und Formeln}
% \begin{definition}{Begriff}
%   ...
% \end{definition}

\section{Verfahren Schritt für Schritt}
% \begin{enumerate}
%   \item ...
% \end{enumerate}

\section{Beispiele}
% \begin{beispiel}{Titel}
%   ...
% \end{beispiel}

\section{Prüfungsrelevant}
% \begin{pruefung}
%   ...
% \end{pruefung}

\section{Offene Fragen}
% Was ist noch unklar? Beim Dozenten oder in der Übung nachfragen.
